\section{Data and measurement}\label{sec:data}

\subsection{The ENAHO quarterly employment modules}

I use the public-use quarterly files of the Encuesta Nacional de Hogares (ENAHO),
the national household survey conducted by Peru's Instituto Nacional de Estad\'istica
e Inform\'atica (INEI). ENAHO is a nationally representative, stratified,
multi-stage survey with continuous fieldwork; the quarterly employment and income
module (module 500) records detailed labor market information for each household
member of working age. I pool the sixteen quarters from the first quarter of 2021
through the fourth quarter of 2024, which yields about 347,000 person-quarter records
of working-age individuals, and I add the four quarters of 2025 for the out-of-sample
validation exercise. All statistics use the module's sampling weights; the treatment of
standard errors is discussed in Section~\ref{sec:strategy}.

I restrict the sample to individuals of working age, defined as ages 14 to 65, and
I drop observations with missing education. My main wage analysis is further
restricted to dependent wage employees, for whom the minimum wage is the relevant
statutory floor, while the employment, participation, formality, and self-employment
outcomes are defined over the working-age population or over the employed, as
appropriate.

\subsection{Education and the definition of skill groups}

The treatment and comparison groups are defined by educational attainment, recorded
in ENAHO variable \texttt{p301a}, the highest level of schooling completed. I
validated the coding of this variable against the weighted distribution of attainment
in the data; it corresponds to the standard INEI classification, ranging from no
schooling through postgraduate study. I define \emph{low-skilled} workers as those
with at most complete secondary schooling (no education, initial, primary, or
secondary), and \emph{high-skilled} workers as those with any tertiary schooling
(non-university higher, university, or postgraduate). I exclude the small residual
category of special education, which is not ordered on the skill dimension, and
observations with missing education.

This partition has three attractive features. First, it is a bright line in the
Peruvian labor market: completion of tertiary schooling is the main credential that
separates workers who compete for jobs paying near the minimum wage from those who do
not. Second, it splits the working-age population into groups of comparable size,
roughly two thirds low-skilled and one third high-skilled, which gives adequate
statistical power on both sides. Third, and most important for identification, the two
groups differ sharply in their exposure to the minimum wage. As
Table~\ref{tab:desc} shows, in the pre-reform period the median monthly wage of a
low-skilled dependent employee is close to the statutory floor, whereas the median
high-skilled employee earns well above it. The classification follows a long line of
work that treats education as a primary determinant of where the minimum wage bites
\citep{lemos_2009, katzkowicz_2021_uruguay}.

\subsection{Outcome variables}

The outcomes fall into three groups. On \emph{pay}, the real hourly wage is monthly labor
income in the main occupation divided by monthly hours, deflated by the Lima consumer price
index to December 2021 soles and winsorized at the first and last half-percentiles; I build
monthly labor income following the ENAHO convention of \citet{jimenez_2023}, and real monthly
earnings is the same numerator without the hours adjustment. On \emph{quantities}, employment
and labor force participation are indicators defined over the working-age population, and
weekly hours is hours worked in the main occupation. On \emph{composition}, formal employment,
self-employment, and an indicator for being paid below the statutory minimum are defined among
the employed; the last is a direct measure of the bite of and compliance with the policy.

A measurement issue arises for informality. INEI's official informal-employment
indicator, \texttt{ocupinf}, is present in the 2021 through 2023 files but is not
released in the 2024 and 2025 files. Because my study window extends through 2024, I
reconstruct a consistent informality indicator from underlying primitives that are
available in every year: pension affiliation, the possession of a formal labor
contract, the registration status of the productive unit, and the category of
employment. I validate the reconstruction against the official \texttt{ocupinf} in
2021 through 2023, where both are available, and obtain agreement of about 88 percent
at the worker level, with an implied informality rate within one percentage point of
the official figure. I use the official indicator wherever it exists and the
reconstruction only for 2024 and 2025, and I report a robustness check that restricts
the formality analysis to the years covered by the official measure. Full details of
the construction are in Appendix~\ref{app:data}.

\subsection{Descriptive statistics}

Table~\ref{tab:desc} reports weighted means of the main variables by skill group in
the pre-reform period. The two groups differ in the expected ways. Low-skilled
workers are older on average, more likely to be male and rural, and have far lower
wages. The employment rate is similar across groups, but the composition of
employment differs starkly: low-skilled workers are much more likely to be informal
and self-employed, and far more likely to earn below the statutory minimum. Roughly
seven in ten employed low-skilled workers earn below the new minimum in the pre-reform
period, against fewer than four in ten high-skilled workers, which is the exposure gap on
which my design rests. The high informality of the low-skilled group,
around 85 percent against about 49 percent for the high-skilled, is central to
interpreting the results: for most low-skilled workers the statutory floor is not
directly enforced, so the margin along which the labor market can adjust is the
allocation of workers across formal and informal arrangements rather than the level of
pay within formal jobs.

\begin{table}[H]\centering
\caption{Descriptive statistics by skill group, pre-reform period}\label{tab:desc}
\begin{threeparttable}
\begin{tabular}{lccc}
\toprule
 & Low-skilled & High-skilled & Difference \\
 & ($\le$ sec. complete) & ($\ge$ higher) & \\
\midrule
Age & 35.92 & 34.65 & 1.27 \\
Female & 0.50 & 0.50 & 0.00 \\
Urban & 0.76 & 0.93 & -0.17 \\
Married & 0.53 & 0.41 & 0.12 \\
Years of education & 8.57 & 14.47 & -5.90 \\
Labour force participation & 0.73 & 0.80 & -0.07 \\
Employed & 0.69 & 0.73 & -0.04 \\
Wage employee (incl. domestic) & 0.29 & 0.48 & -0.19 \\
Self-employed & 0.39 & 0.24 & 0.15 \\
Formal (if employed) & 0.13 & 0.47 & -0.34 \\
Informal (if employed) & 0.87 & 0.53 & 0.34 \\
Weekly hours (if employed) & 39.74 & 41.44 & -1.70 \\
Real hourly wage (S/) & 5.74 & 9.08 & -3.34 \\
Real monthly earnings (S/) & 904.30 & 1,755.12 & -850.82 \\
Below MW (wage earners) & 0.43 & 0.19 & 0.24 \\
\midrule
Observations & 65,954 & 27,278 & \\
\bottomrule
\end{tabular}
\begin{tablenotes}\footnotesize
\item Notes: Weighted means using ENAHO survey weights, working-age individuals (14--65) in the pre-reform quarters (2021Q1--2022Q1). Low-skilled: education at most complete secondary (\texttt{p301a}$\le$6). High-skilled: any tertiary education (\texttt{p301a}$\ge$7). Wage-earner variables (hourly wage, below MW) are conditional on being a private-sector wage earner (employees and domestic workers, monthly-equivalent earnings); hours and formality are conditional on employment. Observation counts are unweighted. The below-MW row uses the statutory floor in force in the income reference month (930 soles throughout the pre-reform window); the exposure shares relative to the NEW floor (57 vs.\ 40 percent) discussed in the text use 1{,}025 soles.
\end{tablenotes}
\end{threeparttable}\end{table}


Figure~\ref{fig:trends} plots the weighted quarterly means of the main outcomes by skill
group over the study window, with the reform marked. The raw series already hint at the
results that follow: the wage series of the two groups move in parallel throughout, while
the formality and self-employment series begin to diverge around the reform. These are
unconditional means; the regression analysis that follows adjusts for composition and
seasonality.

\begin{figure}[H]
\centering
\includegraphics[width=0.95\textwidth]{fig2_raw_trends.pdf}
\caption{Raw outcome trends by skill group, 2021--2024. Weighted quarterly means for
low- and high-skilled workers. The dashed vertical line marks the May 2022 reform.}
\label{fig:trends}
\end{figure}

Figure~\ref{fig:bite} displays the regional variation in the bite of the minimum wage,
measured by the Kaitz index, the ratio of the new minimum to the pre-reform median
dependent wage in each department. The index ranges from about 0.79 in the richest
department to about 2.79 in the poorest, so that in much of the country the new minimum
exceeds the median wage. This variation is predetermined with respect to the reform and
provides the basis for the triple-difference and continuous-intensity checks in
Section~\ref{sec:robustness}.

\begin{figure}[H]
\centering
\includegraphics[width=0.72\textwidth]{fig3_regional_bite.pdf}
\caption{Regional variation in the bite of the minimum wage. The Kaitz index is the
ratio of the new minimum wage (1,025 soles) to the department's median monthly wage of
dependent employees in the pre-reform quarters. Departments above the dashed line have
a median wage below the new minimum.}
\label{fig:bite}
\end{figure}
